\documentclass[11pt]{article}
\usepackage[letterpaper,margin=1in]{geometry}
\usepackage{fontspec}
\setmainfont{texgyretermes}[Extension=.otf,UprightFont=*-regular,BoldFont=*-bold,ItalicFont=*-italic,BoldItalicFont=*-bolditalic]
\usepackage{microtype}
\usepackage{xcolor}
\usepackage[switch]{lineno}
\usepackage{fancyhdr}
\usepackage[numbers,square,sort&compress]{natbib}
\usepackage{enumitem}
\usepackage{placeins}
\renewcommand{\topfraction}{0.9}\renewcommand{\bottomfraction}{0.8}\renewcommand{\textfraction}{0.04}\renewcommand{\floatpagefraction}{0.75}\setcounter{topnumber}{3}\setcounter{totalnumber}{4}
\setlist{itemsep=2pt,topsep=4pt}
\usepackage{graphicx,booktabs,tabularx,array}
\usepackage[labelfont=bf,font=small,labelsep=period,skip=6pt]{caption}
\usepackage{xurl}
\usepackage[hidelinks]{hyperref}

\input{numbers.tex}  % \V{key}: every number comes from build.py (ledger: numbers.csv)
\newcommand{\missing}[1]{\textcolor{red}{[#1]}}

\pagestyle{fancy}\fancyhf{}
\cfoot{\footnotesize\thepage}
\renewcommand{\headrulewidth}{0pt}
\setlength{\parskip}{4pt}

\begin{document}

\begin{center}
{\Large\bfseries Specializing a Decision Model for Radiology: Comparison with General-Purpose Decision Models and Language Models\par}
\vspace{12pt}
Udbhav Ram$^{1,2}$ and Ran Zhang$^{1,2}$\par
\vspace{6pt}
{\small $^{1}$Department of Medical Physics, University of Wisconsin--Madison, Madison, WI, USA\\
$^{2}$Department of Radiology, University of Wisconsin--Madison, Madison, WI, USA\\[4pt]
Correspondence: Ran Zhang, rzhang229@wisc.edu\par}
\end{center}

\vspace{6pt}
\begin{abstract}\setlength{\parskip}{6pt}
\noindent\textbf{Background:} System One-like decision models return probabilities over a provided answer space to pointed questions in a single forward pass, a format well suited to categorical radiology decisions, but a general-purpose decision model has been less accurate than large language models (LLMs) on medical tasks.

\noindent\textbf{Methods:} The general-purpose decision models Kev-27B (built on a Qwen3.8-27B backbone) and Kev-9B were fine-tuned on \V{train_records} public radiology and medical records to produce radiology-specific decision models (RadKev-27B, RadKev-9B) and compared on \V{test_questions} held-out questions with \V{n_generalist} general-purpose decision models and two LLMs, Qwen3.8-27B and MedGemma-27B-text. LLMs, which generated some labels, were compared only on the \V{hk_questions} human-labeled questions. The prespecified primary outcome was accuracy averaged across \V{n_tasks} tasks relative to Kev-27B; other analyses were done post hoc.

\noindent\textbf{Results:} Specialization increased task-averaged accuracy by \V{rb_prim} percentage points (95\% CI \V{rb_prim_ci}), and by \V{rb_prim_hk} points (\V{rb_prim_hk_ci}) over the \V{n_hk_tasks} human-labeled tasks. On human-labeled questions, its gain exceeded that of a threefold increase in model size (\V{rb_spec27} vs \V{rb_scale} points; difference \V{rb_did}, 95\% CI \V{rb_did_ci}), but not when accuracy was averaged over tasks (\V{rb_spec27_tm} vs \V{rb_scale_tm}). The shared backbone was less accurate as a general-purpose decision model than as an LLM (\V{rb_acc_stock27}\% vs \V{rb_acc_qwen38}\%) and more accurate after specialization (\V{rb_acc_v2_27}\%; +\V{rb_rk27_qwen} points, 95\% CI \V{rb_rk27_qwen_ci}), the highest of all systems (MedGemma-27B-text, \V{rb_acc_medgemma_brief}\%). Initializing from Kev-9B rather than its base model added \V{rb_init} points on human-labeled questions (\V{rb_init_ci}). Without the case, the Eurorad gain persisted, indicating cues in the answer options.

\noindent\textbf{Conclusions:} On human-labeled questions, specialization made a general-purpose decision model more accurate than its backbone used as an LLM; its advantage over a threefold increase in size depended on task weighting.
\end{abstract}

% ================================================================ introduction
\section{Introduction}

Many day-to-day radiology workflows involve categorical decision making: whether a report describes a particular finding, which differential diagnosis within a specific subset of options is most likely, to which subspecialty a study should be assigned, and how urgently a result requires triage and follow-up. Increasing automation in such decisions requires a model that selects consistently and confidently among predefined answers and assigns a probability to its selection, so that highly confident decisions can be automated and uncertain ones can be forwarded and flagged for human review \citep{geifman2017selective}. Large language models (LLMs) have been tested in these settings, but they often return free, generated text that must be mapped onto the admissible answers, which can be unreliable and inconsistent, and externally hosted LLMs may offer better performance but are often unsuitable for patient reports because of institutional privacy constraints \citep{mukherjee2023vicuna}.

Decision models, introduced by TypeSafe as so-called System One models, are designed for this setting \citep{almeida2026jev}. A decision model receives an input text, the state, together with typed questions (binary, multiple-choice or ordinal) whose admissible answers are specified in the request body, and returns a probability distribution over the admissible answers to each question. All questions are answered in a single query, in parallel, rather than by autoregressive generation \citep{almeida2026jev,typesafe2026api}, so every output is a valid answer by construction. Jev, released by TypeSafe in September 2026, is available only through an externally hosted interface, so that submitted data must leave the user's institution, and its architecture remains closed source \citep{typesafe2026api,hume2026unmasked}.

Kev is an open-weight family of decision models, with 0.8 to 27 billion parameters, that implements the same TypeSafe-like interface on open-source Qwen backbones \citep{qwen38}, with an architecture inferred from Jev's external behavior \citep{palmer2026kev,hume2026unmasked}. Because its weights and training code are released under the Apache-2.0 license, Kev can be deployed within an institutional network and adapted to domain-specific data. These properties correspond closely to the requirements of radiology: the decisions are categorical and frequent, the several decisions required for one report can be answered in a single request, calibrated probabilities permit selective automation, and patient data need not leave the institution.

General-purpose decision models have, however, been less accurate than LLMs on several medical tasks. In a benchmark evaluation, Jev achieved accuracy comparable to that of the frontier LLM GPT-6 Sol with reasoning on questions derived from research abstracts, but lower accuracy on examination questions and substantially lower accuracy on diagnostic case challenges (59.8\% vs 82.4\% on DiagnosisArena), although its probabilities were better calibrated on examination questions \citep{madrid2026jevmed}. In radiology, Jev has been used to evaluate the factual consistency of generated reports \citep{huang2026jevrad}. For LLMs, the value of domain specialization is contested: domain-adaptive pretraining improves downstream performance \citep{gururangan2020dontstop}, whereas a generalist model with optimized prompting has exceeded medically fine-tuned models \citep{nori2023medprompt}. For decision models, domain adaptation has been reported for Chinese-language medical, legal and financial tasks \citep{wang2026chinesejev}. To our knowledge as of this writing, no decision model has been specialized for radiology, and specialization has not been compared with the principal alternatives for use in radiology clinics: larger general-purpose decision models and LLMs of comparable size.

In this study, we fine-tuned two general-purpose decision models, Kev-27B and Kev-9B, on \V{train_records} records from public radiology and medical sources to obtain RadKev-27B and RadKev-9B, and evaluated them on \V{test_questions} held-out questions under an analysis plan specified before any test result was examined. Specialization was compared with model scale, represented by \V{n_generalist} general-purpose decision models of \V{params_min} to \V{params_max} parameters, and with two LLMs of the same size, Qwen3.8-27B \citep{qwen38} and the medical model MedGemma-27B-text \citep{sellergren2025medgemma}, evaluated zero-shot with reasoning disabled on all questions and enabled on a sample of them. Because Kev-27B is built on Qwen3.8-27B, the design compares identical backbone weights used as an LLM, as a general-purpose decision model and as a radiology-specialized decision model.

% ================================================================ methods
\section{Methods}

\subsection{Data}\label{sec:data}

The data used for this study are intended to cover three distinct decision-making modalities in radiology for which open-source data are available: reading a report, diagnosing a case and applying medical knowledge, with human-reviewed ground-truth labels preferred wherever possible (Table~\ref{tab:data}). Each record pairs a state, i.e. the input text that the model ingests, with one or more typed questions whose admissible answers are given with the question: yes or no, a choice among named options, or an ordinal score. Report reading is represented by chest radiograph reports from IU/Open-i \citep{demnerfushman2016iu}, whose findings were indexed by human coders, posed as yes/no questions on whether a finding is present or the study is normal and as a four-option question on which finding the report describes, and by chest CT reports from CT-RATE \citep{hamamci2026ctrate}, which were posed in the same way against the dataset's classifier labels for 18 abnormalities. Case diagnosis is represented by teaching cases from Eurorad \citep{eurorad,wanglab2025eurorad}, in which the model reads the clinical history and imaging findings and selects the author's final diagnosis from the case's own differential diagnoses, or one of \V{n_sections} subspecialty sections to which the case should be routed. Medical knowledge is represented by four-option examination questions from MedMCQA \citep{pal2022medmcqa} and MedQA \citep{jin2021medqa}, by the medical subjects of MMLU \citep{hendrycks2021mmlu}, by yes, no or maybe questions on research abstracts from PubMedQA \citep{jin2019pubmedqa} and by expert-level questions with up to ten options from MedXpertQA \citep{zuo2025medxpertqa}. Apart from a small set of patient summaries labeled with imaging-guideline topics \citep{yao2025radcases}, decisions about imaging orders, triage and follow-up have no public answer keys; questions of this kind were therefore generated from public reports and clinical indications and labeled by two LLMs (Section~\ref{sec:llmlabels}). Records were divided into training, development and test splits by patient or case (\V{train_records}, \V{dev_records} and \V{test_records} records), with official test sets retained where available and one instruction wording per question kind withheld from training; IU, MMLU, PubMedQA and MedXpertQA were used only for development and testing. Of the \V{test_questions} test questions, \V{hk_questions} have answers assigned by people (human-labeled questions); the remainder are model-labeled, with answers from the CT-RATE classifier or from teacher agreement. Supplementary Note~S1 shows one record from each source, and Supplementary Note~S2 gives the construction rules.

\begin{table}[t]
  \caption{Data sources, answer keys and split sizes. A dash indicates a source not used for training. CT-RATE training records are a random sample of \V{ctrate_cap} reports. Teacher-labeled counts are those of the labels used for training and the main evaluation.}
  \label{tab:data}
  \centering\small
  \input{generated/tab_data.tex}
\end{table}

\FloatBarrier
\subsection{LLM-derived labels}\label{sec:llmlabels}

Decisions about imaging orders, triage and follow-up are routine in radiology, but no large-scale public datasets provide reliable human-generated answer keys for them, and labeling them by hand at the required scale was not feasible within this study. Candidate questions on these decisions, of ten kinds (Supplementary Table~\ref{tab:teacher}), were generated from reports in CheXpert Plus \citep{chambon2024chexpertplus}, ReXGradient-160K \citep{zhang2025rexgradient} and CT-RATE and from clinical indications in these sources and in Eurorad. Two LLMs acted as teachers (Figure~\ref{fig:teacher}): MedGemma-27B-text \citep{sellergren2025medgemma}, a medical model, and Qwen3.8-27B \citep{qwen38}, a general-purpose model. Each read the state and the question with lettered options, with reasoning disabled, and its answer distribution was the softmax of its next-token logits over the option letters. In the run whose labels were used for training, the reasoning segment that MedGemma-27B-text opens by default was not closed, so its letter probabilities were read at the start of that segment rather than after it, which can penalize the model; the labels were later regenerated with the segment closed (Supplementary Note~S2). A question was retained only if both teachers ranked the same answer option first; the agreed-upon answer became the label, and the mean of the two distributions a soft training target \citep{hinton2015distilling}. Of \V{teacher_cand} candidates, \V{teacher_kept_first} (\V{teacher_kept_first_pct}\%) were retained for training. Machine-derived labels are widely used in radiology datasets \citep{irvin2019chexpert,hamamci2026ctrate}, LLM annotation has performed comparably to crowd-worker annotation on text-classification tasks \citep{gilardi2023chatgpt}, and agreement between independent labelers is an established filter for label errors \citep{brodley1999mislabeled}. Because LLM labels still require validation against human judgment \citep{pangakis2023validation}, results on the teacher-labeled test questions are reported as agreement with the teachers (Supplementary Table~\ref{tab:teacher_agreement}), and all conclusions rest on human-labeled questions.

\begin{figure}[t]
  \centering\includegraphics[width=\textwidth]{figures/fig_teacher.pdf}
  \caption{Two-teacher labeling, illustrated with two order questions generated from Eurorad case presentations (cases \V{tex_agreement_case} and \V{tex_disagreement_case}). Each teacher's answer distribution is the softmax of its next-token logits over the option letters. When both teachers rank the same option first (top), the question is retained, with that option as its label and the mean of the two distributions as its soft training target; otherwise (bottom), it is discarded.}
  \label{fig:teacher}
  \par\vspace{2pt}{\footnotesize Case text \textcopyright{} European Society of Radiology, reproduced under CC BY-NC-SA 4.0; this figure is distributed under the same license.\par}
\end{figure}

\FloatBarrier
\subsection{Model and training}\label{sec:training}

\begin{figure}[!t]
  \centering\includegraphics[width=0.86\textwidth]{figures/fig_kev.pdf}
  \caption{Architecture of Kev, the decision model fine-tuned in this study, shown for Kev-27B. (a) A request consists of a state followed by one or more questions, each with its options and a decision token; every question attends to the state and to itself but not to the other questions, so the state is encoded once. (b) The backbone is Qwen3.8-27B: \V{bb_layers} layers in 16 repeating groups of three Gated DeltaNet (linear-attention) layers and one gated full-attention layer, with hidden size \V{bb_hidden}. Its weights are frozen, and a rank-\V{lora_rank} low-rank adapter (LoRA) is attached to every linear projection. (c) The pointer head projects the hidden states of the decision token and of each option's closing token to \V{head_dim} dimensions; the score of an option is the scaled dot product of the two projections, and the answer distribution is the softmax over the question's options after division by a temperature T. Kev-9B has the same structure on Qwen3.5-9B-Base.}
  \label{fig:kev}
\end{figure}

\begin{figure}[!tb]
  \centering\includegraphics[width=\textwidth]{figures/fig_training.pdf}
  \caption{Training. (a) Training loss of RadKev-27B and RadKev-9B, logged every 10 optimizer steps (thin lines) and as a moving average over 500 steps (thick lines); the loss is the cross-entropy over each question's options on the radiology training records and the replayed Kev records. (b) Training loss of the 9B model initialized from Kev-9B (RadKev-9B) and from Qwen3.5-9B-Base with a newly initialized adapter and head (learning rate \V{lr_base}), with otherwise identical data and settings. (c) Accuracy on the development split, by source, of the released Kev models (open markers) and of the fine-tuned models (filled markers); blue, 27B; orange, 9B.}
  \label{fig:training}
\end{figure}

Kev \citep{palmer2026kev} adds a low-rank adapter (LoRA) \citep{hu2022lora} of rank \V{lora_rank} to every linear projection of a frozen language-model backbone, together with a pointer head \citep{vinyals2015pointer} (Figure~\ref{fig:kev}): the state is encoded once, each question attends to the state but not to the other questions, and the resultant answer distribution is a softmax over scores computed between the question and each of its options. RadKev-27B and RadKev-9B were obtained by continued training of the released Kev-27B and Kev-9B, whose backbones are Qwen3.8-27B and Qwen3.5-9B-Base, so that the specialized models keep Kev's interface and differ from it only in their weights. The training split was interspersed with \V{replay} records from Kev's own training data to limit the loss of general decision skill, and Kev's prescribed augmentations were applied: options were reshuffled in each pass, and questions with a single correct answer were randomly modified with ``none of the above'' options or irrelevant distractors, whereas teacher-labeled questions were only reshuffled to preserve their soft targets. The loss was the cross-entropy with respect to the labeled answer or the soft target. Each model was trained for one epoch (\V{steps_27} optimizer steps) with AdamW and a one-cycle schedule, at a peak learning rate of \V{lr_27} for RadKev-27B and \V{lr_9} for RadKev-9B and an effective batch size of \V{accum} records, in bfloat16 on NVIDIA RTX A6000 GPUs (RadKev-27B split across two GPUs, \V{hours_27} h; RadKev-9B on one, \V{hours_9} h). Training loss fell from \V{loss_rk27_start} to \V{loss_rk27_end} for RadKev-27B and from \V{loss_rk9_start} to \V{loss_rk9_end} for RadKev-9B, and development accuracy rose from \V{dev_b27}\% to \V{dev_a27}\% and from \V{dev_b9}\% to \V{dev_a9}\% (Figure~\ref{fig:training}). A temperature was then fitted for each model on the development split by minimizing the negative log-likelihood \citep{guo2017calibration} (\V{T_27} and \V{T_9}). As prespecified, RadKev-27B trained on all sources was retained as the primary model, because its development accuracy was not more than 1.0 percentage point below that of a model trained without the CT-RATE and teacher-labeled data (\V{sel_primary}\% vs \V{sel_ablation}\%). Supplementary Note~S3 gives the model revisions, the remaining hyperparameters and the software versions.

\FloatBarrier
\subsection{Evaluation}\label{sec:evaluation}

The analysis was conducted with a specification of the primary model (RadKev-27B), the primary comparator (Kev-27B), the secondary comparators (Kev-9B, Laya, Qwen3.8-27B and MedGemma-27B-text, and, as ablations, the pilot models and a 27B model trained without the CT-RATE and teacher-labeled data), the outcomes, the model-selection rule and the statistical method; all other systems and analyses are reported as post hoc. Every system was scored on the same \V{test_records} test records (\V{test_questions} questions), and decision models received each record once, with all of its questions in a single request. The comparators separate the effects of specialization, scale and model type: Kev-27B, the starting point of RadKev-27B, measures specialization; Kev-0.8B, Kev-4B and Kev-9B \citep{palmer2026kev} measure scale within the same family; Laya \citep{laya2026}, Laya-typed-decisions \citep{layatyped2026} (421M parameters each), GLiNER2.5-Decide \citep{gliner25decide2026,zaratiana2025gliner2} (486M) and Julia-1 \citep{julia2026} (144M) are the other released general-purpose decision models; Qwen3.8-27B, the backbone of Kev-27B and RadKev-27B, measures the same network queried directly as an LLM; and MedGemma-27B-text is an LLM specialized for medicine. The LLMs were scored with the teacher prompt, one question per prompt, from the softmax of the next-token logits over the option letters with reasoning disabled, and, post hoc, additionally with reasoning enabled on a stratified sample of \V{reason_q} human-labeled questions (Supplementary Note~S4). Because the LLMs had produced the teacher labels, they were compared only on human-labeled questions. The primary outcome was task-averaged accuracy, the unweighted mean of the accuracies of the \V{n_tasks} test tasks, compared between RadKev-27B and Kev-27B after excluding, as prespecified, a demonstration sample of \V{playground_n} records drawn from the test split before the final evaluation; it was also computed over the \V{n_hk_tasks} human-labeled tasks, on which all conclusions are based. Secondary outcomes were accuracy over all questions, by decision family and on withheld instruction wordings; expected calibration error and Brier score; the rate of confident errors; coverage at a 5\% error budget \citep{geifman2017selective}; and general decision performance on Kev's out-of-domain transfer suite (\V{transfer_n} questions).

\FloatBarrier
\subsection{Statistical analysis}\label{sec:stats}

All comparisons were paired: both systems were evaluated on identical questions, and differences are reported as the first system minus the second. Ninety-five percent confidence intervals were obtained by a cluster bootstrap \citep{efron1993bootstrap} with \V{n_boot} resamples, in which test records were resampled within each source together with all of their questions, so that questions about the same report or case were never separated. One set of resamples was shared by all systems, so that every difference, and every difference between two differences, was computed on the same resamples. Intervals are the 2.5th and 97.5th percentiles of the resampled statistic; for task-averaged accuracy, the statistic is the mean of the per-task differences, and differences in calibration, confident errors and coverage were estimated on the first \V{n_boot_cal} resamples. Intervals for Kev's transfer suite, which consists of different questions, were obtained with 1,000 resamples. No formal sample-size calculation was performed; the test set comprised all eligible held-out records of each source (Table~\ref{tab:data}). Two-sided $p$ values were derived from the bootstrap distribution and adjusted by the Holm procedure \citep{holm1979simple} within each set of comparisons reported together; this adjustment was not part of the analysis plan. The demonstration sample was excluded from the primary comparison but not from the other analyses, which use the full test set.

\FloatBarrier
\subsection{Additional analyses}\label{sec:additional}

The following analyses were added after the primary results were known; apart from the initialization ablation, none involved training. The effect of specialization was compared with that of scale as a difference of differences, the gain from specialization (RadKev-27B minus Kev-27B, or RadKev-9B minus Kev-9B) minus the gain from scale (Kev-27B minus Kev-9B), computed on identical resamples. To isolate the contribution of the general-purpose decision model as a starting point, the 9B training was repeated from Qwen3.5-9B-Base with a newly initialized adapter and head, at the default learning rate of Kev's trainer (\V{lr_base}) and otherwise identical data and settings, and both 9B arms were also trained on a fixed 10\% sample of the training split (\V{f10_records} records; \V{f10_steps} optimizer steps). The gain of each specialized model was compared between withheld and seen instruction wordings; calibration was re-estimated after recalibrating every system on the test split by two-fold cross-fitted temperature scaling; and the primary comparison was repeated with the mean over decision families in place of the mean over tasks, alongside a prevalence baseline that always selects the answer position most often correct in each task. An options-only control tested whether the models need the case to answer case-based questions. Every Eurorad diagnosis and MedQA test question was posed again with the case removed: the state, which contains the patient's age, sex, clinical history and imaging findings (Eurorad) or the clinical vignette (MedQA), was replaced by the single sentence ``No case information is available.'', while the instructions and the answer options were left unchanged. A model that relies on the case should then fall to the accuracy expected by chance, one divided by the number of options; accuracy that remains well above chance indicates that the correct option can be recognized from the options themselves. RadKev-27B, Kev-27B, RadKev-9B and Kev-9B were scored in this way, and their accuracy was compared with that on the full case. In addition, Eurorad diagnosis questions were divided according to whether a distinctive word of the correct option, at least six letters long and absent from the other options, appears in the case text. The effect of the answer space was examined on the \V{as_n_dx} Eurorad diagnosis questions and \V{as_n_medqa} MedQA test questions by holding the case and the correct answer fixed and replacing the other options with randomly chosen or semantically similar options from the same dataset, or with plausible incorrect answers generated by Qwen3.8-27B (Supplementary Note~S5). Qwen3.8-27B and MedGemma-27B-text were scored on the same answer spaces with the procedure used for the main evaluation, which labels each option with a single letter and reads the answer from the next-token probabilities of the letters; they were therefore evaluated, and timed, only on answer spaces of at most 16 options. Latency was measured on \V{lat_records} test records (\V{lat_q} questions), one request at a time, on one NVLink pair of RTX A6000 GPUs with Hugging Face Transformers in bfloat16; a decision model answered all questions of a record in one pass, whereas an LLM was called once per question. Two external test sets with expert answer keys, drawn from sources not used in training, were evaluated with the same procedure (Supplementary Note~S6). From RadCases \citep{yao2025radcases}, whose patient summaries are labeled with the American College of Radiology (ACR) Appropriateness Criteria panel and topic that govern their imaging work-up, the \V{ext_rcp_n_all} summaries with openly available text were posed as a choice among the 11 panels and an option stating that no topic applies, and \V{ext_rct_n_all} as a choice among the \V{ext_rc_ntopics} topics; the LLMs were scored on the panel question only, because single-letter scoring is limited to 26 options. From RadGraph-XL \citep{delbrouck2024radgraphxl}, in which radiologists annotated each finding of chest CT, abdominal and pelvic CT and brain MRI reports as present, absent or uncertain, up to three findings per report were posed as the four-option status question used in training (\V{ext_rg_n_all} questions on \V{ext_rg_records} reports).
\FloatBarrier
\section{Results}

\subsection{Primary comparison}\label{sec:res_primary}

\begin{figure}[!t]
  \centering\includegraphics[width=\textwidth]{figures/fig_primary.pdf}
  \caption{Accuracy of Kev-27B (grey) and of RadKev-27B (blue), the model specialized from it, on the full test set. (a) Average over tasks: the unweighted mean of per-task accuracies over all \V{n_tasks} tasks (the primary outcome) and over the \V{n_hk_tasks} human-labeled tasks. (b) Human-labeled decision families, whose answers were set by people (report indexers, case authors and examiners). (c) Model-labeled decision families, whose answers come from the CT-RATE classifier or from two-teacher agreement. Numbers above each pair (\(\Delta\)) are differences in percentage points, with 95\% confidence intervals given in the text; n, test questions per family (\V{test_questions} in total, \V{hk_questions} human-labeled). Excluding the demonstration sample, as prespecified for the primary outcome, changed no difference by more than \V{demo_maxdelta} percentage points. The accuracy axis spans 5 percentage points below the lowest and above the highest value.}
  \label{fig:primary}
\end{figure}

Specialization was found to have improved the prespecified primary outcome, with significant gains in six of the nine decision families: with the demonstration sample excluded, task-averaged accuracy was \V{p_all} percentage points higher for RadKev-27B than for Kev-27B (95\% CI \V{p_all_ci}). On the full test set (Figure~\ref{fig:primary}), it rose from \V{p_all_k27}\% to \V{p_all_rk27}\% (\V{p_all_full} points, \V{p_all_full_ci}), and excluding the demonstration sample changed no difference by more than \V{demo_maxdelta} points and no conclusion. The gain was concentrated in model-labeled tasks: the \V{n_teacher_tasks} teacher-labeled tasks, about a third of the \V{n_tasks}, accounted for \V{share_teacher}\% of it, and over the \V{n_hk_tasks} human-labeled tasks the prespecified difference was \V{p_hk} points (\V{p_hk_ci}). Among the human-labeled families, accuracy improved on chest radiograph report reading (\V{pff_report_cxr_human} points, \V{pff_report_cxr_human_ci}), case diagnosis (\V{pff_case_diagnosis}, \V{pff_case_diagnosis_ci}), a gain that derives partly from the answer options (Section~\ref{sec:res_robust}) and medical knowledge (\V{pff_medical_knowledge}, \V{pff_medical_knowledge_ci}); subspecialty routing did not change detectably (\V{pff_routing}, \V{pff_routing_ci}), and radiology knowledge, with \V{pff_radiology_knowledge_n} questions, did not meet the prespecified criterion (\V{pff_radiology_knowledge}, \V{pff_radiology_knowledge_ci}). Among the model-labeled families, agreement with the model-generated labels increased on CT report reading (\V{pff_report_ct}), imaging orders (\V{pff_orders_protocols}) and triage and follow-up (\V{pff_triage_followup}), but not on chest radiograph finding status (\V{pff_report_cxr}, \V{pff_report_cxr_ci}). All six improvements remained significant after Holm adjustment across the nine families. Averaging over decision families instead of tasks gave a difference of \V{fm_all} points (\V{fm_hk} on human-labeled families), and every system exceeded the prevalence baseline on every human-labeled task (Supplementary Tables~\ref{tab:supp_primary} and~\ref{tab:supp_pertask}). As prespecified, RadKev-27B was also compared with a 27B model trained without the CT-RATE and teacher-labeled data and with the pilot models (Supplementary Table~\ref{tab:supp_ablation}): it exceeded the model trained without these data by \V{abl_v1med_hk} points on human-labeled questions (\V{abl_v1med_hk_ci}) and by \V{abl_v1med_tm} points in task-averaged accuracy (\V{abl_v1med_tm_ci}), so most of the gain on human-labeled questions was obtained without the model-labeled training data.

\FloatBarrier
\subsection{Decision models and language models}\label{sec:res_llm}

\begin{figure}[!t]
  \centering\includegraphics[width=\textwidth]{figures/fig_llm.pdf}
  \caption{Decision models and language models. (a) Accuracy on the \V{hk_questions} human-labeled test questions. (b) Accuracy on the \V{reason_q}-question sample on which the LLMs were also scored with reasoning enabled (Supplementary Note~S4); the decision models answer without reasoning. The accuracy axes in (a) and (b) span 5 percentage points below the lowest and above the highest value. (c) Median latency per question on \V{lat_records} test records, one request at a time on two RTX A6000 GPUs (log scale). A decision model answers all questions about a record in one pass; Qwen3.8-27B was scored from the option-letter logits, by generating the answer letter, or with reasoning (\V{lt_qthink_n} questions).}
  \label{fig:llm}
\end{figure}

On human-labeled questions, RadKev-27B was the most accurate system evaluated (\V{ha_v2_27}\%; Figure~\ref{fig:llm}a). It exceeded Qwen3.8-27B, its own backbone used as an LLM (\V{ha_qwen38}\%; \V{hd_rk27_qwen} points, 95\% CI \V{hd_rk27_qwen_ci}), and MedGemma-27B-text (\V{ha_medgemma_fix}\%; \V{hd_rk27_mg} points, \V{hd_rk27_mg_ci}), whereas the released Kev-27B was less accurate than Qwen3.8-27B by \V{hd_qwen_k27} points (\V{hd_qwen_k27_ci}). RadKev-9B (\V{ha_r9}\%) exceeded Kev-27B by \V{hd_rk9_k27} points (\V{hd_rk9_k27_ci}), and the four other general-purpose decision models reached at most \V{ha_others_max}\% (Supplementary Table~\ref{tab:supp_accuracy}). In a post hoc analysis with reasoning enabled, the accuracy of Qwen3.8-27B on the reasoning sample rose from \V{rs_qwen38}\% to \V{rs_qwen38_think}\% (\V{rp_qthink} points, \V{rp_qthink_ci}) and that of MedGemma-27B-text from \V{rs_medgemma_fix}\% to \V{rs_medgemma_think}\% (Figure~\ref{fig:llm}b). RadKev-27B, without reasoning, reached \V{rs_v2_27}\% on the same questions and did not differ detectably from Qwen3.8-27B with reasoning (\V{rp_rk_qthink} points, \V{rp_rk_qthink_ci}); averaged over tasks, Qwen3.8-27B with reasoning was more accurate (RadKev-27B minus Qwen3.8-27B, \V{rp_rk_qthink_tm} points, \V{rp_rk_qthink_tm_ci}). RadKev-27B was better calibrated than Qwen3.8-27B with reasoning (expected calibration error \V{rs_v2_27_ece} vs \V{rs_qwen38_think_ece}) and could accept more of the sample at a 5\% error budget (\V{rs_v2_27_cov}\% vs \V{rs_qwen38_think_cov}\%). The median latency per question was \V{lt_rk27} ms for RadKev-27B and \V{lt_rk9} ms for RadKev-9B, compared with \V{lt_q} ms for Qwen3.8-27B scored from its option-letter logits, \V{lt_qgen} ms for a generated answer letter and \V{lt_qthink} s with reasoning, about \V{lt_ratio} times the latency of RadKev-27B (Figure~\ref{fig:llm}c); Qwen3.8-27B generated a median of \V{rs_q_medtok} reasoning tokens, and \V{rs_q_cap}\% of its reasoning reached the \V{reason_think_cap}-token limit (Supplementary Tables~\ref{tab:supp_reasoning} and~\ref{tab:supp_latency}).

\FloatBarrier
\subsection{Specialization, scale and initialization}\label{sec:res_spec}

In post hoc comparisons on human-labeled questions, specialization raised accuracy by \V{ss_spec27} percentage points at 27B (95\% CI \V{ss_spec27_ci}) and by \V{ss_spec9} at 9B (\V{ss_spec9_ci}), whereas the threefold increase in size from Kev-9B to Kev-27B raised it by \V{ss_scale} (\V{ss_scale_ci}; Figure~\ref{fig:spec}a). The gain from specialization exceeded the gain from scale by \V{ss_did27} points at 27B (\V{ss_did27_ci}) and by \V{ss_did9} at 9B (\V{ss_did9_ci}), and RadKev-9B was more accurate than Kev-27B (\V{ss_r9k27} points, \V{ss_r9k27_ci}). On the human-labeled radiology questions, scale lowered accuracy (\V{ss_rscale} points, \V{ss_rscale_ci}) while specialization raised it at both sizes (Figure~\ref{fig:spec}b). When accuracy was averaged over tasks rather than questions, however, the gain from scale (\V{ss_tm_scale} points, \V{ss_tm_scale_ci}) exceeded that from specialization at 27B (\V{ss_tm_spec27}, \V{ss_tm_spec27_ci}), so the comparison of specialization with scale depends on how tasks are weighted. In the post hoc initialization ablation, starting from Kev-9B rather than from its base model improved accuracy on human-labeled questions by \V{in_full} points with all of the training data (\V{in_full_ci}) and by \V{in_f10} points with 10\% of it (\V{in_f10_ci}; Figure~\ref{fig:spec}c); with the base-model start, more data did not improve accuracy (\V{sa_b9f10}\% with 10\% and \V{sa_b9}\% with all of the data). Because the base-model start used the default learning rate of Kev's trainer (\V{lr_base}), the two starts also differ in learning rate. Starting from Kev-9B also preserved general decision performance outside radiology (\V{tr_k} points relative to Kev-9B, \V{tr_k_ci}), whereas starting from the base model reduced it (\V{tr_b} points, \V{tr_b_ci}; Figure~\ref{fig:spec}d); for RadKev-27B the change was \V{tr_27} points (\V{tr_27_ci}; Supplementary Table~\ref{tab:supp_init}).

\begin{figure}[!htb]
  \centering\includegraphics[width=\textwidth]{figures/fig_spec.pdf}
  \caption{Specialization, scale and initialization. Top: the released Kev models (grey) and their specialized counterparts, RadKev (blue), at 9B and 27B, on all human-labeled questions (a) and on the human-labeled radiology questions (b; IU, Eurorad and MedMCQA radiology). The numbers above each pair (\(\Delta\)) are the gains from specialization in percentage points; the gain from scale is the difference between the two grey bars. Bottom: the 9B model fine-tuned from Qwen3.5-9B-Base with a new adapter and head (tan) or from Kev-9B (blue), on 10\% or all of the training data. (c) Accuracy on all human-labeled questions; \(\Delta\), Kev-9B start minus base-model start. (d) Change in accuracy on Kev's out-of-domain transfer suite (\V{transfer_n} questions) relative to the released Kev-9B; whiskers, 95\% confidence intervals. Each accuracy axis spans 5 percentage points below the lowest and above the highest value.}
  \label{fig:spec}
\end{figure}


\FloatBarrier
\subsection{Calibration and selective prediction}\label{sec:res_calib}

On human-labeled questions, specialization made more of the questions answerable at a fixed error rate. Answering questions in descending order of confidence, RadKev-27B could answer \V{cal_rk27_cov}\% of them while keeping the error among the answered questions at or below 5\%, compared with \V{cal_k27_cov}\% for Kev-27B (\V{cp_k27_cov5} points, 95\% CI \V{cp_k27_cov5_ci}), \V{cal_q_cov}\% for Qwen3.8-27B (\V{cp_q_cov5} points, \V{cp_q_cov5_ci}) and \V{cal_mg_cov}\% for MedGemma-27B-text (Figure~\ref{fig:calib}a,~b), and its Brier score was lower than that of Kev-27B (\V{cal_rk27_brier} vs \V{cal_k27_brier}). Its probabilities were, however, less well calibrated than those of Kev-27B: the expected calibration error was \V{cal_rk27_ece}, compared with \V{cal_k27_ece} (\V{cp_k27_ece}, \V{cp_k27_ece_ci}), and \V{cal_rk27_ce}\% of questions were answered incorrectly with a confidence of at least 0.9, compared with \V{cal_k27_ce}\% (\V{cp_k27_conf_err} points, \V{cp_k27_conf_err_ci}); the corresponding values were \V{cal_q_ece} and \V{cal_q_ce}\% for Qwen3.8-27B and \V{cal_mg_ece} and \V{cal_mg_ce}\% for MedGemma-27B-text (Figure~\ref{fig:calib}c). After post hoc cross-fitted recalibration on the test split, the calibration errors of RadKev-27B and Kev-27B were equal (\V{rc_rk27_ece} and \V{rc_k27_ece}), and the advantage of RadKev-27B in coverage remained (\V{rc_rk27_cov}\% vs \V{rc_k27_cov}\%; Supplementary Table~\ref{tab:supp_calib}).

\begin{figure}[!htb]
  \centering\includegraphics[width=\textwidth]{figures/fig_calib.pdf}
  \caption{Calibration and selective prediction on the \V{hk_questions} human-labeled test questions. (a) Error among the answered questions when questions are answered in descending order of confidence; the dashed line marks a 5\% error budget. (b) Proportion of questions that can be answered with at most 5\% error. (c) Expected calibration error over ten equal-width confidence bins. In (b) and (c), dark bars show the models as scored (each with its own temperature) and light bars after cross-fitted temperature recalibration on the test split, which leaves accuracy unchanged.}
  \label{fig:calib}
\end{figure}

\FloatBarrier
\subsection{Robustness of the gains}\label{sec:res_robust}

The gain of RadKev-27B over Kev-27B on human-labeled questions was \V{wd_held_out} points on withheld instruction wordings (95\% CI \V{wd_held_out_ci}) and \V{wd_seen} points on wordings seen in training (\V{wd_seen_ci}), a difference of \V{wd_did} points (\V{wd_did_ci}). In the post hoc options-only control, in which the case was removed and only the question and its answer options were shown, RadKev-27B still selected the correct Eurorad diagnosis in \V{bl_dx_rk27}\% of questions, compared with \V{bl_dx_k27}\% for Kev-27B and \V{bl_chance_dx}\% expected by chance (Figure~\ref{fig:robust}a); the gain from specialization without the case (\V{bl_dx_gain_blind} points) thus exceeded the gain with the full case (\V{bl_dx_gain_full}). This result indicates that the case author's final diagnosis could be recognized among the differential diagnoses from the wording of the options alone, and that RadKev-27B relied on this regularity more than Kev-27B. On MedQA, options-only accuracy fell close to chance (25\%) for both models (\V{bl_mq_rk27}\% and \V{bl_mq_k27}\%; Figure~\ref{fig:robust}b; Supplementary Table~\ref{tab:supp_blind}). The Eurorad gain was not explained by words of the correct diagnosis remaining in the case text: it was \V{aw_flag} points (\V{aw_flag_ci}) in the \V{aw_flag_n} questions in which a distinctive word of the correct option occurred in the case and \V{aw_none} points (\V{aw_none_ci}) in the \V{aw_none_n} in which none did (Figure~\ref{fig:robust}d). With the case available, the gain depended on the alternatives offered: it was \V{as_spec27_dx_orig} points with the case's own differential diagnoses, \V{as_spec27_dx_origllm} when LLM-written alternatives were added, \V{as_spec27_dx_llm8} with eight LLM-written options and \V{as_spec27_dx_sim8} with the eight most similar diagnoses from the dataset, but \V{as_spec27_dx_rand64} (\V{as_spec27_dx_rand64_ci}) with 64 randomly chosen diagnoses, for which both models were close to ceiling (Figure~\ref{fig:robust}c). Part of the gain in case diagnosis therefore derives from regularities of the answer options rather than from reading of the case.

\begin{figure}[!htb]
  \centering\includegraphics[width=\textwidth]{figures/fig_robust.pdf}
  \caption{Robustness of the gains. Top: options-only control. Accuracy of Kev-27B (grey) and RadKev-27B (blue) on Eurorad diagnosis (a) and MedQA (b) test questions given the full case, and given only the question and its answer options, with the case removed (replaced by ``No case information is available.''); dashed lines, accuracy expected by chance given the number of options. A model that needs the case should fall to chance when the case is removed. Bottom: the Eurorad diagnosis gain of RadKev-27B over Kev-27B, with 95\% confidence intervals, (c) for different sets of options offered with the correct diagnosis (the case's own differential diagnoses; these with up to \V{as_add} LLM-written alternatives added; eight LLM-written options; the eight most similar or 64 random diagnoses from the dataset; Supplementary Note~S5) and (d) according to whether a distinctive word of the correct option occurs in the case text.}
  \label{fig:robust}
\end{figure}

\FloatBarrier
\subsection{Size and composition of the answer space}\label{sec:res_answerspace}

The accuracy and latency of this new class of decision models depend almost entirely on the complexity, size and completeness of the answer space. Fundamentally, even the most accurate decision models will show poor accuracy and latency benchmarks if the answer space is not correctly defined. In the post hoc answer-space analysis, accuracy fell as alternatives were added to the correct answer, and it depended more on how similar the alternatives were to the correct answer than on how many were offered (Supplementary Tables~\ref{tab:answer_space} and~\ref{tab:answer_space_llm}). With random alternatives from the dataset, the Eurorad accuracy of RadKev-27B fell only from \V{as_rk27_dx_rand2}\% with two options to \V{as_rk27_dx_rand64}\% with 64, whereas with the most similar alternatives it fell from \V{as_rk27_dx_sim2}\% to \V{as_rk27_dx_sim64}\%; for every decision model and both sources, a two-option question with the most similar alternative was answered less accurately than a 64-option question with random alternatives. With the most similar alternatives, RadKev-27B was the most accurate system on Eurorad diagnosis at every size (Figure~\ref{fig:answerspace}a); with eight options, for example, its accuracy was \V{as_rk27_dx_sim8}\%, compared with \V{as_k27_dx_sim8}\% for Kev-27B, \V{asl_q_dx_sim8}\% for Qwen3.8-27B and \V{asl_mg_dx_sim8}\% for MedGemma-27B-text. On MedQA, its advantage was confined to small answer spaces: it led with two options (\V{as_rk27_mq_sim2}\%, against \V{asl_q_mq_sim2}\% for Qwen3.8-27B) but not with 16 (\V{as_rk27_mq_sim16}\%, against \V{asl_q_mq_sim16}\%; Figure~\ref{fig:answerspace}b). The LLMs were evaluated with at most 16 options because of how they are scored: the options are presented as a lettered list and the answer is read from the probability that the model assigns to each option's letter as its next token, so every option must be labeled by a single letter (A to P in the scoring procedure; at most 26 letters in any case). With more than 26 options, the labels would span several tokens and a single next-token readout would no longer give one probability per option. Decision models fundamentally have no such limit, because the pointer head scores each option directly from its own representation, without generating tokens, and Kev accepts up to 255 options per question. Within this range, with the original options, they reproduced the accuracy of the main evaluation for Qwen3.8-27B (\V{asl_q_dx_orig}\% on Eurorad diagnosis) and came within one percentage point of it for MedGemma-27B-text (\V{asl_mg_dx_orig}\% vs \V{asl_mg_orig_main}\%). With one question per request, the latency of RadKev-27B rose from \V{as_lat_rk27_2} ms with two options to \V{as_lat_rk27_16} ms with 16 and \V{as_lat_rk27_128} ms with 128, and that of RadKev-9B from \V{as_lat_rk9_2} to \V{as_lat_rk9_128} ms; the LLMs, scored from their option-letter logits, took a similar \V{asl_lat_q_2} to \V{asl_lat_q_16} ms (Qwen3.8-27B) and \V{asl_lat_mg_2} to \V{asl_lat_mg_16} ms (MedGemma-27B-text) for two to 16 options (Figure~\ref{fig:answerspace}c). The latency curves of the LLMs end at 16 options for the reason given above, not because latency became prohibitive; the decision models were timed with up to 128 options.

\begin{figure}[!htb]
  \centering\includegraphics[width=\textwidth]{figures/fig_answer_space.pdf}
  \caption{Accuracy and latency as the answer space grows. (a, b) Accuracy on \V{as_n_dx} Eurorad diagnosis questions (a) and \V{as_n_medqa} MedQA questions (b) when the correct answer is offered together with the alternatives from the same dataset that are most similar to it; the case is always shown. The LLMs are shown up to 16 options because they were scored from the next-token probability of each option's letter, which requires every option to be labeled by a single letter (A to P); the decision models score each option directly and were evaluated with up to 64 options (a, b) and 128 options (c). (c) Median latency per request with one question per request on \V{as_lat_cases} Eurorad cases. Results with random and LLM-written alternatives, and confidence intervals for the decision models, are given in Supplementary Tables~\ref{tab:answer_space} and~\ref{tab:answer_space_llm}; Supplementary Note~S5 describes the construction of the answer spaces.}
  \label{fig:answerspace}
\end{figure}

\FloatBarrier
\subsection{External test sets}\label{sec:res_external}

On the external test sets (Supplementary Table~\ref{tab:supp_external}), the decision models were more accurate than the LLMs in selecting the ACR Appropriateness Criteria panel, whereas specialization did not improve accuracy at 27B. On the panel question, RadKev-27B reached \V{ext_rcp_v2_27_all}\%, compared with \V{ext_rcp_qwen38_all}\% for Qwen3.8-27B (\V{ext_rcp_d_v2_27_qwen38_all} points, 95\% CI \V{ext_rcp_d_v2_27_qwen38_all_ci}) and \V{ext_rcp_medgemma_fix_all}\% for MedGemma-27B-text, but \V{ext_rcp_stock27_all}\% for Kev-27B (\V{ext_rcp_d_v2_27_stock27_all} points, \V{ext_rcp_d_v2_27_stock27_all_ci}). The difference from Kev-27B arose in the \V{ext_rcp_n_none_only} summaries to which no topic applies, for which RadKev-27B selected the corresponding option less often (\V{ext_rcp_v2_27_none_only}\% vs \V{ext_rcp_stock27_none_only}\%); on the summaries with an applicable panel, RadKev-27B was more accurate (\V{ext_rcp_v2_27_excl_none}\% vs \V{ext_rcp_stock27_excl_none}\%; \V{ext_rcp_d_v2_27_stock27_excl_none} points, \V{ext_rcp_d_v2_27_stock27_excl_none_ci}). At 9B, the difference on the panel question was \V{ext_rcp_d_r9_stock9_all} points (\V{ext_rcp_d_r9_stock9_all_ci}). In topic selection among \V{ext_rc_ntopics} options, RadKev-27B and Kev-27B did not differ (\V{ext_rct_v2_27_all}\% vs \V{ext_rct_stock27_all}\%; \V{ext_rct_d_v2_27_stock27_all} points, \V{ext_rct_d_v2_27_stock27_all_ci}). On the \V{ext_rg_n_all} RadGraph-XL status questions, accuracy was \V{ext_rg_v2_27_all}\% for RadKev-27B and \V{ext_rg_stock27_all}\% for Kev-27B (\V{ext_rg_d_v2_27_stock27_all} points, \V{ext_rg_d_v2_27_stock27_all_ci}), specialization raised it at 9B (\V{ext_rg_d_r9_stock9_all} points, \V{ext_rg_d_r9_stock9_all_ci}), and RadKev-27B was less accurate than Qwen3.8-27B (\V{ext_rg_qwen38_all}\%; \V{ext_rg_d_v2_27_qwen38_all} points, \V{ext_rg_d_v2_27_qwen38_all_ci}) and more accurate than MedGemma-27B-text (\V{ext_rg_medgemma_fix_all}\%).

\FloatBarrier
\section{Discussion}\label{sec:discussion}

Fine-tuning a general-purpose decision model on public radiology and medical data improved the prespecified primary outcome, and on human-labeled questions it produced the most accurate system evaluated. Three findings bear on the choice between specialization and the alternatives. First, on human-labeled questions, specialization raised accuracy more than a threefold increase in model size, and RadKev-9B exceeded Kev-27B; on the human-labeled radiology questions, the larger general-purpose model was less accurate than the smaller one. Second, the same backbone weights were less accurate as a general-purpose decision model than as an LLM and more accurate once the decision model had been specialized, so the general-purpose decision format alone did not account for the gain. Third, the starting point affected the outcome: fine-tuning from Kev-9B rather than from its base model gave higher accuracy, a larger benefit from additional data and no loss of general decision performance. All three findings derive from post hoc analyses and should be regarded as exploratory.

These results extend earlier observations on domain adaptation. For LLMs, continued training on domain text improves downstream performance \citep{gururangan2020dontstop}, whereas prompting alone has allowed a generalist model to exceed medically fine-tuned models \citep{nori2023medprompt}; for clinical and radiology text, lightweight adaptation of open models has produced summaries comparable or superior to those of medical experts \citep{vanveen2024adapted} and is regarded as an alternative to general-purpose models for radiology reports \citep{huisman2025lightweight}, and parameter-efficient fine-tuning of locally deployed open models is recommended where prompting and retrieval are insufficient \citep{bluethgen2025bestpractices,savage2025opensource}. Here, the medically trained LLM, MedGemma-27B-text, was the least accurate of the 27B systems on human-labeled questions. For decision models, an evaluation of Jev found accuracy comparable to that of a frontier LLM on questions about research abstracts but lower accuracy on examination questions and diagnostic cases \citep{madrid2026jevmed}, and domain adaptation has been reported for Chinese-language medical, legal and financial decisions \citep{wang2026chinesejev}. The present results indicate that, for the radiology tasks studied, a specialized open decision model can close the gap to an LLM of the same size and that, at 9B, the general-purpose decision model was a better starting point than its base model, although this comparison is confounded by learning rate. Because no LLM was fine-tuned on the same data, the study does not separate the effect of the decision format from that of domain fine-tuning as such, and the comparison with scale depends on how tasks are weighted: averaged over tasks, the larger general-purpose model gained more than the specialized one, because scale helped more than specialization on examination-style knowledge tasks, such as MedQA and several MMLU subjects, and on subspecialty routing, each of which carries the same weight in the task average as the much larger report-reading tasks.

Compared with LLMs, the specialized model offers a different balance of accuracy, speed and calibration. On the reasoning sample, its accuracy did not differ detectably from that of Qwen3.8-27B with reasoning when every question was weighted equally, although the reasoning LLM was more accurate when accuracy was averaged over tasks, and RadKev-27B answered about \V{lt_ratio} times faster. A decision model also answers all questions about a report in one pass and returns a probability for every admissible answer, which permits automation of confident decisions and referral of uncertain ones \citep{geifman2017selective}: at a 5\% error budget, RadKev-27B could answer \V{cal_rk27_cov}\% of human-labeled questions, more than Kev-27B, Qwen3.8-27B or MedGemma-27B-text. Its probabilities were, however, less well calibrated than those of the model it was specialized from, and it made more confident errors; recalibration on local data, which removed the difference in calibration error in this study, would be required before its probabilities are used as thresholds \citep{guo2017calibration}.

The options-only control qualifies the largest human-labeled gain. On Eurorad diagnosis, RadKev-27B selected the correct diagnosis in most questions without seeing the case, which indicates that the case author's final diagnosis could be recognized among the differential diagnoses from the wording of the options. The gain on this task therefore reflects regularities of the answer options at least as much as reading of the case, and it should not be interpreted as improved diagnostic reasoning. The same caution applies to other benchmarks whose distractors are written by the same authors as the correct answer: a model trained on such data may learn the authors' phrasing rather than the decision.

\subsection{Clinical implementation}\label{sec:disc_clinical}

\paragraph{Suitable decisions.} The considerations in this subsection follow from the present results and from the design of decision models; none of the uses described was evaluated in this study. Of the roles proposed for language models in radiology, which also include report drafting, summarization and patient communication \citep{paschali2025foundation,bluethgen2025bestpractices}, a decision model addresses the categorical decisions and complements, rather than replaces, generative systems. The approach appears suited to frequent, categorical decisions made on radiology text for which the admissible answers can be listed in advance: protocol and examination selection from an order's clinical indication, prioritization of requests and worklists, routing of studies to subspecialty readers, detection of critical or incidental findings in reports, extraction of follow-up recommendations for tracking, and consistency checks between a report's findings and its impression. A decision model answers all questions about a report in one pass and returns a probability for each admissible answer, so a single request can populate several fields of a worklist or a quality dashboard, and confident decisions can be automated while uncertain ones are referred for review. It is less suited to decisions whose answers cannot be listed in advance, to the generation of text such as report impressions, and to decisions that depend on the images themselves.

\paragraph{Resources and hardware.} All components are openly available and can be run on institutional hardware, so patient data need not leave the institution; the RadKev weights are licensed for non-commercial research use only. In this study, the 27B models ran on two GPUs with 48 GB of memory each and the 9B models on one; specialization required \V{hours_27} h of training for RadKev-27B on two such GPUs and \V{hours_9} h for RadKev-9B on one, so retraining on an institution's own data would be feasible on a single workstation with two such GPUs. The 9B model offers a lower hardware requirement at a moderate loss of accuracy on human-labeled questions (\V{ha_r9}\% vs \V{ha_v2_27}\%). Beyond hardware, deployment would require integration with the radiology information system to supply report text and receive decisions, a labeled sample of local cases on which to measure accuracy and refit the temperature, radiologist time to define and review the answer spaces, and monitoring of accuracy and calibration over time.

\paragraph{Latency and throughput.} RadKev-27B answered a question in a median of \V{lt_rk27} ms and RadKev-9B in \V{lt_rk9} ms, one request at a time, compared with \V{lt_qthink} s for an LLM with reasoning. The advantage over LLMs lies in avoiding generation and reasoning and in answering all questions about a report in one pass: an LLM scored only from its option-letter logits was about as fast per question (\V{lt_q} ms), but this form of use returns one answer per request, is limited to answer spaces that single letters can label, and was less accurate on human-labeled questions. At these latencies, one GPU pair answers about \V{thr_rk27_hr} questions per hour without batching (\V{thr_rk9_hr} for the 9B model on one GPU), estimates that are extrapolated from the measured latencies, assume uninterrupted sequential requests and were not measured in a clinical system. Because every option is encoded as part of its question, latency is nearly constant up to about 16 options and increases approximately in proportion to the number of options beyond 32 (Figure~\ref{fig:answerspace}c), which favors short answer spaces.

\paragraph{Choosing the answer space.} A decision model can only choose among the answers it is offered, so the answer space is part of every decision, and the analysis of Section~\ref{sec:res_answerspace} indicates that its composition affected accuracy more than its size. Questions became difficult when the alternatives resembled the correct answer, not when there were many of them, and the advantage of the specialized model was largest when the alternatives resembled the correct answer. Conversely, an answer space that omits the correct answer forces an error, and one filled with implausible alternatives overstates accuracy. For decisions with a closed set of answers, such as the examination to order, the protocol, the priority or the follow-up category, the answer space can be taken from the institution's own order catalog, protocol list, triage scale or structured-reporting template; such lists are defined once, reviewed by radiologists, versioned and audited. For open-ended decisions such as differential diagnosis, a candidate list can be generated for each case and the decision model used to choose among the candidates and to attach a calibrated probability to the choice: candidates may be retrieved from a curated catalog of diagnoses by similarity to the case text, restricted by modality, body region, age and sex, or written by an LLM asked for a differential diagnosis. In this study, alternatives written by an LLM resembled the correct answer about as closely as the case authors' own differential diagnoses did (mean embedding cosine \V{as_cos_llm8} and \V{as_cos_orig}), which suggests that LLM-generated differentials may provide candidates of comparable difficulty; their clinical plausibility was not assessed. Lists of tens rather than hundreds of options keep latency low, and an explicit ``none of the above'' or ``other'' option, which Kev's training procedure supports, provides a means of indicating that the correct answer is missing, although this was not evaluated here.

\paragraph{Safeguards.} Whatever the method of building answer spaces, its inclusion rate, the proportion of locally labeled cases in which the correct answer appears among the candidates, should be measured before deployment, because a decision model cannot recover an answer it is not offered. The temperature should be refitted on local cases and on the answer spaces that will be used, because calibration depends on both, and thresholds for automation should be set on that local sample rather than taken from this study. Decisions with clinical consequences should remain subject to review by a radiologist, as the model's developers also recommend \citep{palmer2026kev}, and accuracy and calibration should be monitored after deployment, because case mix, report style and order catalogs change over time.

\subsection{Limitations}\label{sec:limitations}

The questions on imaging orders, triage and follow-up were labeled by agreement of two LLMs without validation against human judgment \citep{pangakis2023validation}; their indications, when taken from Eurorad, sometimes named the examination performed, which can reveal the answer to order questions, and results on these tasks therefore measure agreement with the teachers rather than accuracy. Most human-labeled radiology questions came from a single collection of chest radiograph reports (IU/Open-i), and on these questions the margin of RadKev-27B over Qwen3.8-27B was small (\V{ss_rad_rk27_q} points, \V{ss_rad_rk27_q_ci}); its larger margin on all human-labeled questions derived mainly from the medical examination questions. On two external test sets that were not part of training, specialization did not improve accuracy at 27B (Section~\ref{sec:res_external}). All data were drawn from public datasets and teaching collections, the models were not validated on institutional reports, and accuracy was not examined across patient subgroups or institutions. The LLMs were evaluated zero-shot with a single prompt, and no LLM was fine-tuned on the same data. MedGemma-27B-text was evaluated without reasoning in the main evaluation, with a prompt revised after its first test results had been read, and its accuracy on MedQA (\V{bl_mq_mg_full}\%) was below that reported in its technical report \citep{sellergren2025medgemma}; its results may therefore underestimate the capability of the model. The comparison of starting points differs in learning rate as well as in initialization, and every model was trained once, with a single random seed. Public examination questions and Eurorad cases may have been part of the pretraining data of every model, so paired differences, rather than absolute accuracies, are the basis of the conclusions. The analysis plan was not publicly registered: its commit is held in a private repository, although the plan is reproduced in the code repository. The demonstration sample was excluded only from the primary comparison, and all analyses other than the primary comparison and the prespecified secondary comparisons were added after the primary results were known and are exploratory. Latency was measured one request at a time, and that of the LLM with reasoning on only \V{lt_qthink_n} questions with a shorter token limit than in the accuracy evaluation; LLMs served with batching may achieve higher throughput. Finally, all questions were posed on text, and the study does not address decisions that require the images themselves.

\subsection{Future work}\label{sec:future}

Several questions remain open. The labels derived from LLM agreement should be validated against radiologists' judgments, and the specialized model should be compared with radiologists in a reader study on institutional reports, in which the answer spaces are defined by the institution's own catalogs. A language model fine-tuned on the same data would separate the effect of the decision format from that of domain fine-tuning, and repeated training with different random seeds would quantify the variability of the gains. Methods for constructing answer spaces, by retrieval from curated catalogs or by LLM-generated differentials, should be compared by their inclusion rate and by the downstream accuracy and calibration they yield. Training with alternatives drawn from different sources than the correct answer, for example LLM-written alternatives, may prevent the model from learning to recognize the correct option by its wording. Finally, extending the approach to decisions made on images as well as on text would address the larger part of radiological practice.

\section{Conclusions}\label{sec:conclusions}

Specializing an open, general-purpose decision model on public radiology and medical data produced a model that, on human-labeled questions, was more accurate than its own backbone used as an LLM and than a medically trained LLM of the same size. In post hoc analyses, specialization improved accuracy more than a threefold increase in model size when every question was weighted equally, but not when accuracy was averaged over tasks, and fine-tuning from the general-purpose decision model was more effective than fine-tuning from its base model. On a sample of human-labeled questions, the accuracy of the specialized model did not differ detectably from that of an LLM with reasoning when every question was weighted equally, at about one hundredth of the latency, and it could answer the largest share of questions within a fixed error budget. Two qualifications apply: part of its gain on case diagnosis derives from regularities of the answer options rather than from the case, and its probabilities require recalibration on local data before they are used to decide which cases to automate. For categorical radiology decisions that must run within an institution, specializing a general-purpose decision model may be a practical alternative to scaling it up or to querying an LLM, subject to validation on local data.

\section*{Declarations}

\paragraph{Data availability.} All data used in this study are publicly available. IU/Open-i \citep{demnerfushman2016iu} is distributed under CC BY-NC-ND 4.0, Eurorad \citep{eurorad,wanglab2025eurorad} under CC BY-NC-SA 4.0, MedMCQA \citep{pal2022medmcqa} under Apache-2.0, MedQA \citep{jin2021medqa} under CC BY 4.0, and MMLU \citep{hendrycks2021mmlu}, PubMedQA \citep{jin2019pubmedqa} and MedXpertQA \citep{zuo2025medxpertqa} under MIT licenses; CT-RATE \citep{hamamci2026ctrate} (CC BY-NC-SA 4.0), CheXpert Plus \citep{chambon2024chexpertplus} and ReXGradient-160K \citep{zhang2025rexgradient} are available from their providers under access agreements. The RadCases labels \citep{yao2025radcases} are released under the MIT license with openly available case texts \citep{chen2024medbullets}, and RadGraph-XL \citep{delbrouck2024radgraphxl} is available from the Stanford Center for Artificial Intelligence in Medicine and Imaging under its research use agreement. The question records built from Eurorad, MedMCQA, MedQA, MMLU, PubMedQA and MedXpertQA are released in a Hugging Face dataset (\url{https://huggingface.co/datasets/ramu9703/radkev-data}), each under the license of its source. Because the license of IU/Open-i does not permit derivative works and the access-controlled sources may not be redistributed, records built from them are not released as text; instead, the record identifiers (report, case or patient), SHA-256 hashes of the record text, the answer keys, the teacher labels and a script that rebuilds the records from the original datasets are provided. Per-question model outputs (probabilities for every answer option, without text) are released for all systems and analyses.

\paragraph{Code and model availability.} The code for data construction, teacher labeling, training, evaluation and statistical analysis, and the script that generates every number, table and figure in this manuscript from the job outputs, are available at \url{https://github.com/udiram/RadKev}. The RadKev-27B and RadKev-9B adapters and pointer heads are released on Hugging Face (\url{https://huggingface.co/ramu9703/radkev-27b-v2} and \url{https://huggingface.co/ramu9703/radkev-9b}) under CC BY-NC-SA 4.0 for non-commercial research use, because several training sources carry that license; access requires acceptance of these terms. The Kev code on which they depend is distributed under Apache-2.0 \citep{palmer2026kev}.

\paragraph{Ethics statement.} This study used only publicly available, de-identified datasets under their respective licenses and access agreements. No patients were recruited and no institutional patient data were used; review by an institutional review board was therefore not required.

\paragraph{Author contributions.} U.R. conceived the study, designed the methodology, wrote the software, built the datasets, trained and evaluated the models, performed the statistical analysis, prepared the figures and wrote the manuscript. R.Z. supervised the study and reviewed and edited the manuscript. Both authors approved the final manuscript.

\paragraph{Funding.} This work was supported by the Departments of Radiology and Medical Physics, University of Wisconsin--Madison.

\paragraph{Competing interests.} The authors declare no competing interests.

\paragraph{Use of AI tools.} An AI assistant (Claude, Anthropic) was used under the authors' direction to write and run analysis code, to generate figures and to draft and revise the manuscript text. Every number in the manuscript is generated by code from the job outputs, and the authors reviewed all code, results and text and take full responsibility for the content.

\paragraph{Acknowledgments.} Training and evaluation were run on a workstation of the Zhang laboratory with four NVIDIA RTX A6000 GPUs. We thank the creators of IU/Open-i, Eurorad, CT-RATE, CheXpert Plus, ReXGradient-160K, MedMCQA, MedQA, MMLU, PubMedQA and MedXpertQA for making their data available, the European Society of Radiology and the authors of the Eurorad teaching cases, and the developers of Kev, Qwen and MedGemma for releasing their models openly.

\FloatBarrier
\bibliographystyle{unsrtnat}
\bibliography{refs}

% ================================================================ supplementary information
\clearpage
\section*{Supplementary Information}
\setcounter{table}{0}\renewcommand{\thetable}{S\arabic{table}}
\setcounter{figure}{0}\renewcommand{\thefigure}{S\arabic{figure}}

\subsection*{Supplementary Note S1. Example records}
For each open-license source, the record shown is the one at the first quartile of state length among the \V{playground_n} demonstration records described in Section~\ref{sec:evaluation} (60 per source), reproduced as the models received it. The correct option is shown in bold. For CT-RATE and the teacher-labeled questions, whose report text cannot be redistributed, the question formats are shown instead.

\input{generated/examples.tex}

\FloatBarrier
\subsection*{Supplementary Note S2. Record construction}

\paragraph{IU/Open-i.} Medical Subject Headings assigned by the indexers were mapped to \V{iu_vocab} findings. A finding was considered present if one of its terms had been assigned and absent otherwise, and a study was considered normal if it had been indexed as normal. Reports with fewer than 20 characters of findings and impression, and reports without indexing, were excluded. Each report received two presence questions, chosen so that one finding was present and one absent where possible, one normal-study question and, for half of the reports with at least one finding, a question on which of four findings (one present, three absent) was described. The license (CC BY-NC-ND) does not permit derivative works, so the collection was used only for development (25\% of reports, by a hash of the report identifier) and testing (75\%).

\paragraph{CT-RATE.} The collection comprises 25,692 chest CT studies from 21,304 patients \citep{hamamci2026ctrate}; its labels were produced by a RadBERT classifier fine-tuned on manually annotated reports. One report per study received three presence questions and the normal-study and which-finding questions; a study was normal if none of the 18 abnormalities was labeled. The official training set was divided by patient into training and development records (8\% development), the official validation set served as the test set, and training used a random sample of \V{ctrate_cap} reports.

\paragraph{Eurorad.} The state contained age, sex, clinical history and imaging findings, from which sentences mentioning pathology, biopsy, histology, surgery, treatment, follow-up, confirmation or diagnosis were removed by keyword matching. Cases whose final diagnosis did not appear among the differential diagnoses, that had fewer than two differential diagnoses or whose remaining findings were shorter than 40 characters were excluded, as were development and test cases whose remaining text contained the diagnosis. Half of the cases received the routing question, whose options were shown with a one-line description of the section in 60\% of cases. Age and sex were taken from the case metadata and occasionally differ from those stated in the history, as in the example of Supplementary Note~S1. Cases were split 80\%, 10\% and 10\% by a hash of the case identifier.

\paragraph{Examination questions.} MedMCQA: all radiology questions were retained; other subjects were sampled at approximately 20\% and capped at \V{medmcqa_cap} training questions per subject; questions with empty or duplicate options were removed; because the official test answers are not public, the official validation set served as the test set, and 5\% of the official training set was held out for development. MedQA: the official test set was retained and 10\% of the official training set was held out for development. MMLU: the anatomy, clinical knowledge, college biology, college medicine, medical genetics and professional medicine subjects, with the official validation and test sets used for development and testing. PubMedQA: the expert-annotated subset, divided equally between development and testing. MedXpertQA: the text subset, 20\% for development and 80\% for testing.

\paragraph{Teacher-labeled questions.} Up to \V{teacher_per_source} reports and \V{teacher_per_source} indications were drawn per source, including Eurorad case presentations for indications. Five kinds concern a report (critical finding, urgency, follow-up, interval change and incidental finding), one the status of a finding in a chest radiograph report, and four an imaging order (first examination, use of contrast, priority and appropriateness of an ordered examination); order questions were posed on the clinical indication alone. Each teacher was instructed to answer as an attending radiologist, with reasoning disabled and the options as a lettered list. Because MedGemma's chat template does not suppress its reasoning segment, a one-sentence segment (``I will answer with the letter only.'') was pre-filled so that the next token is the answer. In the run whose labels were used for training, agreement ranged from \V{agree_first_min}\% (use of contrast) to \V{agree_first_max}\% (finding status). The labels used for training, model selection and the main evaluation come from a first run in which MedGemma was scored without the pre-filled segment (\V{teacher_kept_first} questions retained, \V{teacher_kept_first_pct}\%). After the main evaluation, all candidates were rescored by MedGemma with the pre-filled segment and the labels regenerated, with the answers of Qwen3.8-27B unchanged and without retraining. Of the \V{teacher_train_first} training questions retained in the first run, \V{teacher_train_same} (\V{teacher_train_same_pct}\%) were also retained in the regenerated run, and no question retained in both runs received a different label. Supplementary Table~\ref{tab:teacher_agreement} uses the regenerated labels; all other results use the labels of the first run.

\begin{table}[!h]
  \caption{Teacher-labeled question kinds. Candidates were answered by MedGemma-27B-text and Qwen3.8-27B; a question was retained when both ranked the same answer first. Order questions saw only the clinical indication. First run, labels used for training, model selection and the main evaluation; regenerated, labels after rescoring MedGemma with the pre-filled reasoning segment.}
  \label{tab:teacher}
  \centering\footnotesize\setlength{\tabcolsep}{3.5pt}
  \input{generated/tab_teacher.tex}
\end{table}

\FloatBarrier
\subsection*{Supplementary Note S3. Training details}

\paragraph{Augmentation.} Each multiple-choice question with a single correct answer received at most one modification: the correct option was replaced by a ``none of the above'' option, which then became the correct answer (probability \V{p_none}; questions with more than two options); a ``none of the above'' option was added as a wrong alternative (probability \V{p_none_distract}); or an irrelevant distractor was added (probability \V{p_distract}). With probability \V{p_none_pair}, a record additionally contributed a minimal pair in which the same question was posed with a ``none of the above'' option, once with and once without its correct option. With the \V{replay} replayed Kev records, training comprised \V{requested_records} records and \V{examples_seen} examples, including the minimal pairs.

\paragraph{Optimization.} AdamW with weight decay \V{weight_decay}; one-cycle schedule with 10\% warm-up and cosine decay; gradient-norm clipping at 1.0; batch size one with gradient accumulation over \V{accum} records; gradient checkpointing; states limited to \V{max_state} tokens, a limit that no training record exceeded; random seed 0. The temperature was selected from \V{grid_points} logarithmically spaced values between 0.25 and 4. The development comparison for model selection averaged accuracy over the \V{sel_tasks} tasks shared by both development sets.

\paragraph{Models.} Kev-27B, Hugging Face revision \V{rev_kev27}, on Qwen3.8-27B revision \V{rev_qwen27}; Kev-9B, revision \V{rev_kev9}, on Qwen3.5-9B-Base revision \V{rev_qwen9}. The pointer head projects the hidden states of each option's closing token and of the question's decision token into a \V{head_dim}-dimensional space, and the logit of an option is the scaled dot product of the two projections.

\paragraph{Hardware and software.} Peak memory was \V{peak_27} GiB per GPU for RadKev-27B and \V{peak_9} GiB for RadKev-9B. For Kev-9B, splitting the model by layers across two GPUs reproduced single-GPU probabilities exactly on \V{parity_q} development questions. Software: Kev at commit \V{kev_commit}, PyTorch \V{ver_torch} with CUDA \V{ver_cuda}, Transformers \V{ver_transformers}, PEFT \V{ver_peft}, flash-linear-attention \V{ver_fla} and causal-conv1d \V{ver_cc1d}.

\FloatBarrier
\subsection*{Supplementary Note S4. Evaluation details}

\paragraph{Decision models.} Each record was submitted once, with all of its questions, in bfloat16; each model used its own temperature (the released temperature for the general-purpose models and the fitted temperature for RadKev). The answer to a question is the option with the highest probability, and that probability is its confidence. Every model accepted every test question.

\paragraph{Language models.} The prompt was that used for teacher labeling: the state, the question and its options as a lettered list, followed by an instruction to answer with one letter. MedGemma-27B-text was scored with the one-sentence pre-filled reasoning segment described in Supplementary Note~S2. With reasoning enabled, Qwen3.8-27B was run in its thinking mode, and for MedGemma-27B-text the prompt opened its reasoning segment so that the model always reasoned, using vLLM \V{ver_vllm} in bfloat16, each model's default sampling settings, seed 0 and at most \V{reason_think_cap} reasoning tokens. After the reasoning, ``Answer:'' was appended, and the answer distribution was read from the next-token probabilities of the option letters. When the model had written an answer after its reasoning (the first option letter in the generated text) and that answer differed from the letter read-out, the written answer was taken, with a probability of 0.98; this applied to none of the questions for Qwen3.8-27B and to \V{ovr_mg_pct}\% for MedGemma-27B-text. The reasoning sample comprised \V{reason_q} questions in \V{reason_records} records: every Eurorad and MedMCQA radiology question, \V{reason_iu_cap} questions from each IU task and \V{reason_cap} from each other human-labeled task, selected in a fixed hash order that does not depend on any model output.

\paragraph{Outcomes.} The demonstration sample of \V{playground_n} records had been scored for a demonstration interface before the final evaluation; matched by state, it corresponds to \V{demo_records} test records and \V{demo_questions} questions. The expected calibration error was computed over ten equal-width bins of confidence. A confident error is a question answered incorrectly with a confidence of at least 0.9. Coverage at a 5\% error budget is the largest proportion of questions that can be accepted in descending order of confidence while the error among the accepted questions does not exceed 5\%.


\FloatBarrier
\subsection*{Supplementary Note S5. Answer-space analysis}

For each question, the state, the instructions and the human-assigned key were held fixed, and only the other options were varied. Pool distractors were drawn from all distinct options of the same source across the training, development and test splits, either at random or in order of decreasing similarity to the key, measured as the cosine between PubMedBERT sentence embeddings \citep{gu2021pubmedbert}. LLM distractors were generated by Qwen3.8-27B, which received the state, the question and the correct answer and was asked, with reasoning disabled and greedy decoding, for \V{as_n_alt} plausible but incorrect answers; the key was never taken from the LLM. A candidate distractor was discarded if its normalized text contained, or was contained in, that of the key or of an already selected distractor, or if its embedding cosine to either was at least \V{as_dup}. Each question was posed with its original options; with the key and random or most similar pool distractors, for \V{as_sizes} options in total; with the key and LLM distractors, for \V{as_llm_sizes} options; and with the original options supplemented with up to \V{as_add} LLM distractors. An answer space was omitted for a question when too few distractors remained after filtering (16-option LLM answer space: \V{as_n_llm16} of \V{as_q} questions retained). Each answer space was posed as a separate question on the same state. RadKev-27B, Kev-27B, RadKev-9B and Kev-9B were scored on every answer space, and differences were estimated with a paired bootstrap over questions (\V{n_boot} resamples). Latency as a function of the number of options was measured for RadKev-27B and RadKev-9B with one question per request and \V{as_lat_sizes} options (the key and random pool distractors) on \V{as_lat_cases} Eurorad cases, in random order, with the first \V{as_lat_warmup} requests excluded. The LLMs were scored and timed in the same way as in the main evaluation, one question per request, on the answer spaces with at most 16 options, the number of single letters (A to P) used to label options in the scoring procedure; beyond 26 options, single-letter labels are exhausted and a next-token readout no longer yields one probability per option, whereas the decision models score each option directly. In the main latency measurement, the first \V{lat_warmup} requests of each mode (the first request for the LLM with reasoning) were excluded, and the per-question latency of a decision model was its latency per record divided by the number of questions.


\FloatBarrier
\subsection*{Supplementary Note S6. External test sets}

\paragraph{RadCases.} The RadCases release (Hugging Face \texttt{michaelsyao/RadCases}, revision 93c99559) labels each patient summary, a one-sentence description of the presentation, with the ACR Appropriateness Criteria panel and topic that apply to it; the labels were assigned by two senior medical students under the supervision of an attending radiologist, with disagreements resolved by radiologists \citep{yao2025radcases}. The release contains only SHA-512 hashes of the summaries. Their text was rebuilt from the two openly available subsets, summaries written by GPT-3.5 for the original study and the first sentences of Medbullets USMLE Step 2 and 3 questions \citep{chen2024medbullets}, with the sentence splitter of the original study, and a case was retained only if the hash of its rebuilt text matched a label (\V{ext_rc_matched} of \V{ext_rc_labels}). Summaries whose text occurs in any record of the training, development or test splits of this study (\V{ext_rc_overlap}) were excluded, as were, for the panel question, cases labeled with more than one panel (\V{ext_rc_multi}) and, for the topic question, cases labeled with more than one topic or with none. The panel question offered the 11 panels and the option ``None: no ACR Appropriateness Criteria topic applies'' under neutral, shuffled keys; \V{ext_rcp_n_none_only} of the \V{ext_rcp_n_all} panel questions carry this answer. The topic question offered all \V{ext_rc_ntopics} topics of the criteria version used by the original study. The state was the summary alone.

\paragraph{RadGraph-XL.} The Stanford portion of RadGraph-XL comprises \V{ext_rg_reports} reports, 500 each of chest radiographs, chest CT, abdominal and pelvic CT and brain MRI, in which board-certified radiologists annotated every entity and its status \citep{delbrouck2024radgraphxl}. Reports sharing at least half of their word 8-grams with a CheXpert Plus report or with any record of this study were excluded (\V{ext_rg_excl} reports, all of them chest radiograph reports, which derive from the same institution as CheXpert Plus), so that the test consists almost entirely of CT and MRI reports. For each remaining report, one observation annotated as definitely present, one as definitely absent and one as uncertain were drawn at random where available, excluding spans that consist only of a generic qualifier such as ``normal'' or ``stable'', and each was posed in the format of the finding-status questions used in training, with the annotated span as the finding and the options reported as present, explicitly reported as absent, possible or equivocal, and not mentioned in the report; the last was never correct. The state was the full report text.

\paragraph{Analysis.} Accuracy and paired differences were estimated as in the main evaluation, with 2,000 bootstrap resamples of records stratified by subset (RadCases) or modality (RadGraph-XL). Supplementary Table~\ref{tab:supp_external} gives the accuracy of every system.

\FloatBarrier
\subsection*{Supplementary Tables}

\begin{table}[!h]
  \caption{Agreement of the decision models with the regenerated teacher labels of the test split (\V{teacher_test_v2} questions), in percent. These values are agreement with the answer on which both teachers agreed, not accuracy against a human key. RadKev was trained on labels from the first run, which share the procedure and the Qwen3.8-27B answers with the regenerated labels; its agreement is therefore expected to exceed that of Kev and does not indicate correctness.}
  \label{tab:teacher_agreement}
  \centering\footnotesize\setlength{\tabcolsep}{4pt}
  \input{generated/tab_teacher_agreement.tex}
\end{table}

\begin{table}[!h]
  \caption{Primary comparison by decision family: RadKev-27B minus Kev-27B, accuracy in percentage points (95\% confidence interval). Prespecified: with the \V{demo_records} demonstration records excluded, as in the analysis plan; Holm $p$: adjusted across the nine families (not part of the plan). Full test set: all records. The last two rows average the per-family differences instead of the per-task differences (aggregation sensitivity; point estimates).}
  \label{tab:supp_primary}
  \centering\footnotesize\setlength{\tabcolsep}{4pt}
  \input{generated/supp_primary.tex}
\end{table}

\begin{table}[!h]
  \caption{Accuracy (\%) on every test task, full test set. Baseline: always choosing the answer position that is most often correct in the task (prevalence baseline). Qwen, Qwen3.8-27B; MedGemma, MedGemma-27B-text. Top: human-labeled tasks; middle: CT-RATE (classifier labels); bottom: teacher-labeled tasks (first-run labels, as in the main evaluation). LLMs were compared only on human-labeled tasks.}
  \label{tab:supp_pertask}
  \centering\scriptsize\setlength{\tabcolsep}{4pt}
  \input{generated/supp_pertask.tex}
\end{table}

\begin{table}[!h]
  \caption{Accuracy (\%, 95\% confidence interval) of every system. Human-labeled and radiology human-labeled: pooled over questions; all tasks, task mean: the prespecified aggregation. Machine-labeled questions are included in the last two columns for the decision models only; the LLM rows are given for completeness and are not compared with the decision models on them.}
  \label{tab:supp_accuracy}
  \centering\scriptsize\setlength{\tabcolsep}{4pt}
  \input{generated/supp_accuracy.tex}
\end{table}

\begin{table}[!h]
  \caption{Results on the \V{reason_q}-question reasoning sample. Accuracy pooled over questions and averaged over tasks (\%), accuracy on its radiology questions (\%), expected calibration error (ECE) and coverage at a 5\% error budget (\%). Tokens: median number of reasoning tokens and percentage of outputs that reached the \V{reason_think_cap}-token limit.}
  \label{tab:supp_reasoning}
  \centering\footnotesize\setlength{\tabcolsep}{4pt}
  \input{generated/supp_reasoning.tex}
\end{table}

\begin{table}[!h]
  \caption{Latency per question on \V{lat_records} test records (one request at a time, two RTX A6000 GPUs, Hugging Face Transformers, bfloat16; first \V{lat_warmup} requests of each mode excluded, first request for reasoning). For decision models, n counts records and the latency per question is the latency per record divided by its number of questions; Qwen3.8-27B with reasoning was timed on a subsample with a 2,048-token limit.}
  \label{tab:supp_latency}
  \centering\footnotesize\setlength{\tabcolsep}{4pt}
  \input{generated/supp_latency.tex}
\end{table}

\begin{table}[!h]
  \caption{The 9B models: accuracy (\%, 95\% confidence interval) on human-labeled questions and on the human-labeled radiology questions, and change in accuracy on Kev's out-of-domain transfer suite (\V{transfer_n} questions) relative to the released Kev-9B (percentage points, 95\% confidence interval). Models started from Qwen3.5-9B-Base were trained at a learning rate of \V{lr_base}.}
  \label{tab:supp_init}
  \centering\scriptsize\setlength{\tabcolsep}{4pt}
  \input{generated/supp_init.tex}
\end{table}

\begin{table}[!h]
  \caption{Calibration and selective prediction on human-labeled questions: expected calibration error (ECE) over ten equal-width bins, Brier score, confident errors (percentage of questions answered incorrectly with confidence of at least 0.9) and coverage at a 5\% error budget (\%), as scored and after cross-fitted temperature recalibration on the test split (recal.).}
  \label{tab:supp_calib}
  \centering\scriptsize\setlength{\tabcolsep}{4pt}
  \input{generated/supp_calib.tex}
\end{table}

\begin{table}[!h]
  \caption{Options-only control: accuracy (\%, 95\% confidence interval) with the full case and with the case removed and only the question and its options shown.}
  \label{tab:supp_blind}
  \centering\footnotesize\setlength{\tabcolsep}{4pt}
  \input{generated/supp_blind.tex}
\end{table}

\begin{table}[!h]
  \caption{Accuracy (\%) under each answer space. RK, RadKev. Eurorad diagnosis, n = \V{as_n_dx}; MedQA, n = \V{as_n_medqa}. K, number of options offered. Conditions for which too few distractors remained after filtering are evaluated on fewer questions (LLM, K = 16: \V{as_n_llm16} of \V{as_q}). Pool distractors closest to the key may include acceptable alternatives to it that the duplicate filter did not remove; accuracy with similar pool distractors is therefore a lower bound.}
  \label{tab:answer_space}
  \centering\footnotesize\setlength{\tabcolsep}{4pt}
  \input{generated/tab_answer_space.tex}
\end{table}

\begin{table}[!h]
  \caption{Accuracy (\%) of the LLMs under each answer space with at most 16 options (lettered options A--P). Same questions and answer spaces as Supplementary Table~\ref{tab:answer_space}.}
  \label{tab:answer_space_llm}
  \centering\footnotesize
  \input{generated/supp_answer_space_llm.tex}
\end{table}

\begin{table}[!h]
  \caption{Prespecified ablation and pilot comparators, full test set: accuracy (\%, 95\% confidence interval) on human-labeled questions and averaged over all 29 tasks. Differences in percentage points: RadKev-27B minus the model for the 27B models; the pilot minus Kev-9B for the 9B models. The 27B model without CT-RATE and teacher-labeled data was trained on the remaining sources with the recipe of RadKev-27B; the pilot models were trained on the openly licensed sources only, for 400 (27B) and 1,353 (9B) optimizer steps, with 4,000 replayed Kev records.}
  \label{tab:supp_ablation}
  \centering\scriptsize\setlength{\tabcolsep}{4pt}
  \input{generated/supp_ablation.tex}
\end{table}


\begin{table}[!h]
  \caption{Accuracy (\%) on the external test sets. RadCases panel: choice among 11 ACR Appropriateness Criteria panels and ``no topic applies''; excluding None: the questions whose answer is a panel. RadCases topic: choice among 225 topics, decision models only. RadGraph-XL status: four-option status of a radiologist-annotated finding, by modality.}
  \label{tab:supp_external}
  \centering\scriptsize\setlength{\tabcolsep}{4pt}
  \input{generated/supp_external.tex}
\end{table}

\end{document}
